\documentclass[
    12pt,
    openright,
    oneside,
    a4paper,
    english,
    french,
    spanish,
    brazil
]{abntex2}

% ---
% Pacotes básicos
% ---
\renewcommand{\ABNTEXchapterfont}{\bfseries\rmfamily}
\usepackage[T1]{fontenc}
\usepackage[utf8]{inputenc}
\usepackage{lastpage}
\usepackage{indentfirst}
\usepackage{color}
\usepackage{graphicx}
\usepackage{microtype}
\usepackage{caption}
\usepackage{subcaption} 

% ---
% Pacotes para Listagens de Código
% ---
\usepackage{listings}
\renewcommand{\lstlistingname}{Listagem}
\renewcommand{\lstlistlistingname}{Lista de listagens}
\lstset{
    basicstyle=\ttfamily\small,
    breaklines=true,
    frame=single,
    showstringspaces=false
}

% ---
% Pacote de citações (Citações nominais padrão ABNT)
% ---
\usepackage[alf]{abntex2cite}

% ---
% Informações de dados para CAPA e FOLHA DE ROSTO
% ---
\titulo{Relatório Técnico: [Insira o Título do Trabalho Aqui]}
\autor{Samir Soares de Melo}
\local{Montes Claros - MG}
\data{[Ano]}
\instituicao{
    Universidade Estadual de Montes Claros -- Unimontes
    \par
    Centro de Ciências Exatas e Tecnológicas -- CCET
    \par
    Curso de Bacharelado em Sistemas de Informação
}
\tipotrabalho{Relatório Técnico}
\preambulo{Relatório técnico apresentado à disciplina de [Nome da Disciplina], ministrada pelo Prof. [Nome do Professor], referente à atividade [Nome da Atividade].}

% ---
% Configurações de aparência do PDF final
% ---
\definecolor{blue}{RGB}{41,5,195}
\makeatletter
\hypersetup{
    pdftitle={\@title},
    pdfauthor={\@author},
    pdfsubject={\imprimirpreambulo},
    pdfcreator={LaTeX with abnTeX2},
    colorlinks=true,
    linkcolor=black,
    citecolor=black,
    urlcolor=blue,
    bookmarksdepth=4
}
\makeatother

\setlength{\parindent}{1.3cm}
\setlength{\parskip}{0.2cm}

% ----
% Início do documento
% ----
\begin{document}
\selectlanguage{brazil}
\frenchspacing

% ----------------------------------------------------------
% ELEMENTOS PRÉ-TEXTUAIS
% ----------------------------------------------------------
\imprimircapa
\imprimirfolhaderosto

% --- Listas e Sumário ---
\pdfbookmark[0]{\listfigurename}{lof}
\listoffigures*
\cleardoublepage

\pdfbookmark[0]{\lstlistlistingname}{lol}
\lstlistoflistings*
\cleardoublepage

\pdfbookmark[0]{\contentsname}{toc}
\tableofcontents*
\cleardoublepage

% ----------------------------------------------------------
% ELEMENTOS TEXTUAIS
% ----------------------------------------------------------
\textual

% ---
\chapter{Introdução}
% ---
\section{Problema Tratado}
[Descreva aqui o contexto e o problema que o seu relatório ou sistema visa resolver. Lembre-se de escrever sempre na terceira pessoa impessoal, ex: "Observou-se que...", "Nota-se a necessidade de..."]. 

\section{Justificativa}
A elaboração deste trabalho justifica-se pela importância de aplicar os conceitos teóricos na prática. Para o desenvolvimento, utilizou-se o software Exemplo\footnote{Exemplo: Descrição oficial da ferramenta ou linguagem. Disponível em: https://www.link-oficial.com/. Acesso em: [Data].}, o que permitiu consolidar o aprendizado estrutural da disciplina.

\section{Objetivo}
O presente trabalho tem como objetivo [inserir verbo no infinitivo: desenvolver, analisar, documentar] um [inserir o que foi feito], visando [inserir o propósito final].

% ---
\chapter{Fundamentação Teórica}
% ---
Neste capítulo, apresentam-se os conceitos teóricos que embasam a prática. Segundo \citeonline{[chave_da_referencia]}, a teoria é fundamental para a construção de soluções sólidas. 

Além dos artigos científicos, protocolos e padrões também determinam as diretrizes de desenvolvimento. Um exemplo clássico é o protocolo XPTO\footnote{XPTO: Protocolo de exemplo. Disponível em: https://... Acesso em: [Data].}, que define as regras de comunicação do sistema avaliado \cite{[chave_da_referencia2]}.

% ---
\chapter{Desenvolvimento}
% ---
Para garantir a organização e modularidade do trabalho, o projeto foi dividido em etapas lógicas, detalhadas nas subseções a seguir.

\section{[Nome da Primeira Etapa]}
Iniciou-se a implementação configurando a base do projeto. Conforme ilustrado na Listagem \ref{lst:codigo_exemplo}, o algoritmo foi estruturado para receber os dados iniciais.

\begin{lstlisting}[caption={Exemplo de código fonte inicial}, label={lst:codigo_exemplo}, language=PHP]
// Insira seu codigo aqui
$variavel = "Exemplo";
\end{lstlisting}

Dessa forma, a rotina apresentada na listagem garante que os dados sejam processados corretamente antes de avançar para a próxima fase.

\section{[Nome da Segunda Etapa]}
A interface e a organização visual dos elementos foram estruturadas de forma responsiva. Na Figura \ref{fig:exemplo_interface}, observa-se a tela inicial desenvolvida.

\begin{figure}[htb]
    \centering
    \caption{Tela representativa do sistema}
    \label{fig:exemplo_interface}
    % \includegraphics[width=0.8\textwidth]{nome_da_imagem.png}
    \fbox{\rule{0pt}{2in} \rule{0.8\textwidth}{0pt}} % Placeholder (apague e descomente a linha de cima)
    \legend{Fonte: Elaborado pelo autor ([Ano]).}
\end{figure}

A partir da Figura \ref{fig:exemplo_interface}, constata-se que a disposição dos botões facilita a usabilidade do operador. 

Caso o sistema necessite de comparações visuais, podem ser utilizadas figuras compostas. A Figura \ref{fig:exemplo_composto} mostra o estado inicial, enquanto a segunda parte exibe o resultado final.

\begin{figure}[htb]
    \centering
    \begin{subfigure}[b]{0.45\textwidth}
        \centering
        \fbox{\rule{0pt}{1.5in} \rule{1in}{0pt}} % Placeholder
        \caption{Estado Inicial}
        \label{fig:composta_a}
    \end{subfigure}
    \hfill
    \begin{subfigure}[b]{0.45\textwidth}
        \centering
        \fbox{\rule{0pt}{1.5in} \rule{1in}{0pt}} % Placeholder
        \caption{Estado Final}
        \label{fig:composta_b}
    \end{subfigure}
    \caption{Comparação entre diferentes estados do sistema}
    \label{fig:exemplo_composto}
    \legend{Fonte: Elaborado pelo autor ([Ano]).}
\end{figure}

Para não sobrecarregar o texto principal do relatório técnico, as capturas adicionais de tela foram alocadas no Apêndice A.

% ---
\chapter{Trabalhos Correlatos}
% ---
Ao analisar o estado da arte referente ao tema abordado, constata-se que soluções modernas seguem padrões arquiteturais específicos. \citeonline{[chave_estado_da_arte]} elaboraram um estudo onde demonstram que as abordagens atuais reduzem o tempo de processamento em X\%.

Nesse sentido, debates em fóruns especializados complementam a literatura acadêmica. Ao avaliar discussões recentes de desenvolvedores no \textit{Stack Overflow} e em repositórios do \textit{GitHub}, observam-se as seguintes opiniões recorrentes:
\begin{itemize}
    \item A comunidade recomenda a utilização da abordagem Y por facilitar a manutenção a longo prazo.
    \item Profissionais da área apontam que a biblioteca Z pode apresentar vulnerabilidades se não configurada corretamente.
\end{itemize}

Com base nesses apontamentos práticos e acadêmicos, a escolha técnica do presente trabalho procurou alinhar-se com as melhores práticas documentadas.

% ---
\chapter{Conclusão}
% ---
O desenvolvimento deste projeto permitiu a aplicação prática dos conceitos abordados na disciplina. Conclui-se que a elaboração da solução exigiu um certo nível intermediário de planejamento e a adoção de boas práticas de estruturação de código. O sistema atendeu a todos os requisitos propostos, consolidando a proficiência nas ferramentas e tecnologias estudadas.

% ----------------------------------------------------------
% ELEMENTOS PÓS-TEXTUAIS
% ----------------------------------------------------------
\postextual

% Certifique-se de ter o arquivo "referencias.bib" criado no Overleaf
\bibliography{referencias} 

% ---
% Apêndices
% ---
\begin{apendicesenv}
\partapendices

\chapter{[Título do Apêndice: Ex: Telas Adicionais ou Diagramas]}
\label{apendice:telas}

Este apêndice documenta artefatos complementares gerados ao longo do desenvolvimento, evitando a poluição visual do documento principal. A Figura \ref{fig:apendice_extra} demonstra um fluxo de exceção do sistema.

\begin{figure}[htb]
    \centering
    \caption{Fluxo de exceção ou tela adicional}
    \label{fig:apendice_extra}
    % \includegraphics[width=0.9\textwidth]{imagem_extra.png}
    \fbox{\rule{0pt}{2in} \rule{0.9\textwidth}{0pt}} % Placeholder
    \legend{Fonte: Elaborado pelo autor ([Ano]).}
\end{figure}

Dessa forma, conclui-se a documentação visual complementar, atestando o funcionamento do escopo proposto.

\end{apendicesenv}

\end{document}